\chapter{Część teoretyczna}

\section{Taksonomia problemów jakości danych tabelarycznych}
Współczesna analiza danych tabelarycznych oraz uczenie maszynowe opierają się na założeniu, że zbiór uczący w sposób wierny odzwierciedla rozkład badanego zjawiska w populacji generalnej. W warunkach rzeczywistych dane wejściowe są jednak obarczone licznymi zniekształceniami. Jak wskazuje Walker \cite{Walker2021}, błędy te można podzielić na trzy fundamentalne kategorie: braki danych, zaburzenia rozkładów wywołane obserwacjami odstającymi oraz niespójności skali i reprezentacji kategorycznej.

\subsection{Mechanizmy powstawania braków danych}
W teorii statystyki, sformułowanej pierwotnie przez Donalda Rubina \cite{LittleRubin2002}, wyróżnia się trzy podstawowe mechanizmy generowania braków danych:
\begin{enumerate}
    \item \textbf{MCAR (ang. \textit{Missing Completely at Random})} – braki całkowicie losowe. Prawdopodobieństwo wystąpienia braku w danej zmiennej nie zależy ani od wartości samej zmiennej, ani od wartości jakichkolwiek innych atrybutów w zbiorze:
    \begin{equation}
    P(M \mid Y_{\text{obs}}, Y_{\text{mis}}) = P(M),
    \end{equation}
    gdzie $M$ jest binarną macierzą wskaźnikową braków, $Y_{\text{obs}}$ oznacza wartości zaobserwowane, a $Y_{\text{mis}}$ wartości brakujące.
    \item \textbf{MAR (ang. \textit{Missing at Random})} – braki losowe warunkowo. Prawdopodobieństwo braku zależy od innych zaobserwowanych zmiennych, lecz nie od samej wartości brakującej:
    \begin{equation}
    P(M \mid Y_{\text{obs}}, Y_{\text{mis}}) = P(M \mid Y_{\text{obs}}).
    \end{equation}
    Przykładem jest sytuacja, w której częstość braków w wynikach badań laboratoryjnych zależy od wieku pacjenta lub płci.
    \item \textbf{MNAR (ang. \textit{Missing Not at Random})} – braki nielosowe. Prawdopodobieństwo wystąpienia braku zależy bezpośrednio od niezaobserwowanej wartości (np. pacjenci z bardzo wysokim poziomem bólu częściej odmawiają udziału w badaniu).
\end{enumerate}
Wybór odpowiedniej procedury imputacji musi uwzględniać powyższy charakter zjawiska, aby nie doprowadzić do zniekształcenia kowariancji cech \cite{LittleRubin2002,Walker2021}.

\subsection{Wartości odstające i błędy grube}
Wartości odstające (ang. \textit{outliers}) to obserwacje, które w sposób drastyczny odbiegają od wzorca wyznaczonego przez przeważającą część zbioru. Mogą one wynikać z:
\begin{itemize}
    \item \textbf{Błędów grubych aparatury i zapisu}: pomyłki jednostek (np. gramy zamiast miligramów), błędy wprowadzania danych z klawiatury (np. przesunięcie przecinka dziesiętnego), usterki czujników.
    \item \textbf{Naturalnej zmienności ekstremalnej}: rzadkie, lecz rzeczywiste przypadki kliniczne (np. pacjenci z ekstremalnym nasileniem wskaźników stanu zapalnego).
\end{itemize}
Obecność obserwacji odstających silnie zaburza nieodporne estymatory statystyczne. Przykładowo, średnia arytmetyczna $\hat{\mu}$ oraz wariancja próbkowa $\hat{\sigma}^2$ charakteryzują się zerowym punktem załamania (ang. \textit{breakdown point} $\epsilon^* = 0$, wg teorii statystyki odpornej Hubera \cite{Huber1981}), co oznacza, że pojedyncza skrajna wartość może dowolnie zmienić ich oszacowanie. Z kolei estymatory odporne, takie jak mediana czy rozstęp międzykwartylowy, zachowują stabilność przy obecności do $50\%$ zaburzonych obserwacji \cite{Huber1981,Walker2021}.

\section{Metody wstępnego przetwarzania i czyszczenia danych}

\subsection{Uzupełnianie braków danych (Imputacja)}
W praktyce inżynierii danych wyróżnia się podejścia proste (statystyczne) oraz zaawansowane (uczące się):
\begin{enumerate}
    \item \textbf{Imputacja statystyczna}: zastąpienie braków wartością syntetyczną wyznaczoną globalnie na zbiorze uczącym:
    \begin{itemize}
        \item dla cech ciągłych: mediana $\mathrm{Med}(X)$ próby (bezpieczniejsza od średniej przy rozkładach skośnych) \cite{Walker2021},
        \item dla cech kategorycznych: dominanta (moda), czyli najczęściej występująca kategoria.
    \end{itemize}
    \item \textbf{Imputacja sąsiedztwa (\texttt{KNNImputer})}: wielowymiarowa metoda estymacji braków bazująca na algorytmie $k$-najbliższych sąsiadów \cite{SklearnKNNImputer}. Dla dwóch obserwacji $x, y$, w których część cech może brakować, odległość euklidesowa z maskowaniem wyznaczana jest ze wzoru:
    \begin{equation}
    d(x, y) = \sqrt{ \frac{p}{p_{\text{wspólne}}} \sum_{i \in \text{wspólne}} (x_i - y_i)^2 },
    \end{equation}
    gdzie $p$ to całkowita liczba cech, a $p_{\text{wspólne}}$ to liczba cech obserwowanych w obu rekordach jednocześnie. Brakująca wartość w rekordzie $x$ jest następnie estymowana jako średnia ważona z $k$ najbardziej podobnych rekordów:
    \begin{equation}
    \hat{x}_j = \frac{\sum_{m \in N_k(x)} w_m \cdot y_{m,j}}{\sum_{m \in N_k(x)} w_m}, \quad w_m = \frac{1}{d(x, y_m)}.
    \end{equation}
\end{enumerate}

\subsection{Identyfikacja i obsługa wartości odstających}
W pracy zaimplementowano regułę rozstępu międzykwartylowego (IQR) zaproponowaną przez Johna Tukeya \cite{Tukey1977}. Kwartyle $Q_1$ i $Q_3$ dzielą uporządkowany ciąg danych na cztery równe części, a rozstęp międzykwartylowy definiowany jest jako:
\begin{equation}
\mathrm{IQR} = Q_3 - Q_1.
\end{equation}
W klasycznym podejściu Tukeya \cite{Tukey1977} wyróżnia się:
\begin{itemize}
    \item anomalie umiarkowane: $x \notin [Q_1 - 1{,}5 \cdot \mathrm{IQR}, \; Q_3 + 1{,}5 \cdot \mathrm{IQR}]$,
    \item anomalie ekstremalne: $x \notin [Q_1 - 3{,}0 \cdot \mathrm{IQR}, \; Q_3 + 3{,}0 \cdot \mathrm{IQR}]$.
\end{itemize}
Zgodnie z recepturami Walkera \cite{Walker2021}, w niniejszej pracy zastosowano próg $3{,}0 \cdot \mathrm{IQR}$, aby odfiltrować wyłącznie bezdyskusyjne błędy grube i anomalie skrajne. Zamiast agresywnego usuwania całych rekordów (co prowadzi do utraty cennych informacji w pozostałych poprawnych kolumnach), zaimplementowany transformator \texttt{OutlierToNaNTransformer} zamienia wykryte outliery na wartości \texttt{NaN}, umożliwiając ich późniejszą poprawną rekonstrukcję w module imputacji.

Do metod wielowymiarowej detekcji anomalii należy algorytm \textbf{Isolation Forest} \cite{Liu2008}. Metoda ta buduje zespół losowych drzew decyzyjnych (\textit{iTrees}). Punkty odstające, jako rzadkie i odizolowane, wymagają znacznie mniejszej liczby podziałów do ich wyodrębnienia. Wskaźnik anomalii dla obserwacji $x$ definiowany jest jako:
\begin{equation}
s(x, n) = 2^{-\frac{\mathbb{E}(h(x))}{c(n)}},
\end{equation}
gdzie $h(x)$ to długość ścieżki w drzewie, $\mathbb{E}(h(x))$ to średnia długość ścieżki w zespole, a $c(n) = 2\ln(n - 1) + 0{,}5772 - \frac{2(n - 1)}{n}$ jest średnią długością ścieżki nieudanego wyszukiwania w drzewie BST dla $n$ obserwacji. Wartość $s \to 1$ jednoznacznie identyfikuje anomalię \cite{Liu2008}.

\subsection{Skalowanie atrybutów i kodowanie zmiennych jakościowych}
\begin{enumerate}
    \item \textbf{Standaryzacja $z$-score}: przekształca zmienną losową do postaci o zerowej średniej i jednostkowej wariancji:
    \begin{equation}
    z = \frac{x - \mu}{\sigma},
    \end{equation}
    gdzie parametry $\mu, \sigma$ są wyznaczane wyłącznie na próbie uczącej.
    \item \textbf{Skalowanie min-max}: mapuje wartości do zamkniętego przedziału $[0, 1]$:
    \begin{equation}
    x' = \frac{x - x_{\min}}{x_{\max} - x_{\min}}.
    \end{equation}
    \item \textbf{Kodowanie One-Hot (\texttt{OneHotEncoder})}: przekształca cechę kategoryczną o $C$ unikalnych poziomach w $C$ binarnych wektorów wskaźnikowych. W celu uniknięcia problemu rzadkich kategorii wprowadzono transformator \texttt{RareCategoryToNaNTransformer}, który kategorie o częstości występowania poniżej progu 5\% konwertuje do postaci braków danych przed etapem kodowania.
\end{enumerate}

\section{Wybrane algorytmy klasyfikacji}

\subsection{Naiwny Klasyfikator Bayesa (Gaussian Naive Bayes)}
Naiwny Klasyfikator Bayesa opiera się na bezpośrednim zastosowaniu twierdzenia Bayesa przy założeniu warunkowej niezależności cech $x_1, \dots, x_p$ względem zmiennej decyzyjnej $y \in \{0, 1\}$ \cite{Hastie2009,SklearnGaussianNB}:
\begin{equation}
P(y = c \mid x) = \frac{P(y = c) \prod_{j=1}^{p} P(x_j \mid y = c)}{P(x)}.
\end{equation}
W wariancie \texttt{GaussianNB} zakłada się, że rozkład warunkowy każdej cechy ciągłej jest rozkładem normalnym:
\begin{equation}
P(x_j \mid y = c) = \frac{1}{\sqrt{2\pi \sigma_{jc}^2}} \exp\left( -\frac{(x_j - \mu_{jc})^2}{2\sigma_{jc}^2} \right).
\end{equation}
Parametry $\mu_{jc}$ i $\sigma_{jc}^2$ estymowane są za pomocą średniej i wariancji próbkowej. W konsekwencji obecność nieodfiltrowanych wartości odstających drastycznie zniekształca parametry gęstości prawdopodobieństwa, powodując załamanie zdolności predykcyjnej klasyfikatora.

\subsection{Perceptron Wielowarstwowy (MLP)}
MLP (ang. \textit{Multilayer Perceptron}) to sieć neuronowa o architekturze jednokierunkowej, uczona z nadzorem \cite{Hastie2009,SklearnMLPClassifier}. Dla pojedynczego neuronu w warstwie ukrytej sygnał wyjściowy wyznaczany jest jako:
\begin{equation}
a = f\left( \sum_{j=1}^{p} w_j x_j + b \right),
\end{equation}
gdzie $w_j$ oznaczają wagi połączeń, $b$ to obciążenie (bias), a $f(\cdot)$ jest nieliniową funkcją aktywacji (np. ReLU: $f(s) = \max(0, s)$).

Proces uczenia polega na minimalizacji funkcji straty (entropii krzyżowej dla klasyfikacji binarnej):
\begin{equation}
\mathcal{L}(w, b) = -\frac{1}{N} \sum_{i=1}^{N} \left[ y_i \ln \hat{y}_i + (1 - y_i) \ln (1 - \hat{y}_i) \right] + \frac{\alpha}{2} \|w\|_2^2,
\end{equation}
za pomocą algorytmu optymalizacji opartego na gradientach (Adam \cite{KingmaBa2015}). Brak standaryzacji cech wejściowych prowadzi do złego uwarunkowania macierzy drugich pochodnych (hesjanu), co wywołuje zjawisko oscylacji gradientu i uniemożliwia zbieżność procesu uczenia.

\subsection{Lasy Losowe (Random Forest)}
Algorytm Lasów Losowych, opracowany przez Leo Breimana \cite{Breiman2001}, stanowi zespół wielu drzew decyzyjnych uczonych metodą agregacji bootstrapowej (bagging). Każde drzewo $T_b$ budowane jest na niezależnej próbie losowej ze zwracaniem $D_b \subset D$. W każdym węźle drzewa optymalny podział wyznaczany jest wyłącznie spośród losowo wybranego podzbioru $m \approx \sqrt{p}$ cech.

Kryterium oceny jakości podziału węzła stanowi zanieczyszczenie Giniego (ang. \textit{Gini Impurity}):
\begin{equation}
I_G(t) = 1 - \sum_{k=1}^{K} p_k^2,
\end{equation}
gdzie $p_k$ to frakcja obserwacji klasy $k$ w węźle $t$. Ponieważ decyzje podziału w drzewach decyzyjnych zależą wyłącznie od porządku relacyjnego wartości atrybutu ($x_j \le \theta$), Lasy Losowe są całkowicie niewrażliwe na monotoniczne transformacje i brak standaryzacji cech \cite{Breiman2001,Hastie2009}.

\subsection{XGBoost (Extreme Gradient Boosting)}
XGBoost \cite{ChenGuestrin2016} jest algorytmem sekwencyjnego wzmacniania gradientowego. W odróżnieniu od baggingu, każde kolejne drzewo $f_t(x)$ jest dopasowywane do residuów (błędów) popełnionych przez dotychczasowy zespół:
\begin{equation}
\hat{y}_i^{(t)} = \hat{y}_i^{(t-1)} + f_t(x_i).
\end{equation}
Funkcja celu w kroku $t$ przyjmuje postać:
\begin{equation}
\mathrm{Obj}^{(t)} = \sum_{i=1}^{N} l(y_i, \hat{y}_i^{(t-1)} + f_t(x_i)) + \Omega(f_t),
\end{equation}
gdzie człon regularyzacyjny drzewa o $T$ liściach i wagach $w$ definiowany jest jako:
\begin{equation}
\Omega(f) = \gamma T + \frac{1}{2} \lambda \sum_{j=1}^{T} w_j^2 + \alpha \sum_{j=1}^{T} |w_j|.
\end{equation}
Stosując rozwinięcie funkcji straty w szereg Taylora drugiego rzędu, otrzymuje się uproszczoną postać celu:
\begin{equation}
\widetilde{\mathrm{Obj}}^{(t)} \approx \sum_{i=1}^{N} \left[ g_i f_t(x_i) + \frac{1}{2} h_i f_t^2(x_i) \right] + \Omega(f_t),
\end{equation}
gdzie $g_i = \partial_{\hat{y}^{(t-1)}} l(y_i, \hat{y}^{(t-1)})$ oraz $h_i = \partial^2_{\hat{y}^{(t-1)}} l(y_i, \hat{y}^{(t-1)})$. Optymalna waga liścia $j$ wynosi:
\begin{equation}
w_j^* = -\frac{\sum_{i \in I_j} g_i}{\sum_{i \in I_j} h_i + \lambda},
\end{equation}
a zysk z podziału węzła na podzbiory $I_L$ i $I_R$ dany jest wzorem:
\begin{equation}
\mathrm{Gain} = \frac{1}{2} \left[ \frac{(\sum_{i \in I_L} g_i)^2}{\sum_{i \in I_L} h_i + \lambda} + \frac{(\sum_{i \in I_R} g_i)^2}{\sum_{i \in I_R} h_i + \lambda} - \frac{(\sum_{i \in I} g_i)^2}{\sum_{i \in I} h_i + \lambda} \right] - \gamma.
\end{equation}

\section{Metryki oceny jakości modeli i aparat statystyczny}

\subsection{Krzywa ROC i pole pod krzywą (AUC)}
Krzywa ROC przedstawia zależność pomiędzy czułością ($\mathrm{TPR}$) a odsetkiem fałszywych alarmów ($\mathrm{FPR}$) dla zmiennego progu klasyfikacji $\theta \in [0, 1]$ \cite{Fawcett2006}:
\begin{equation}
\mathrm{TPR}(\theta) = \frac{\mathrm{TP}(\theta)}{\mathrm{TP}(\theta) + \mathrm{FN}(\theta)}, \qquad \mathrm{FPR}(\theta) = \frac{\mathrm{FP}(\theta)}{\mathrm{FP}(\theta) + \mathrm{TN}(\theta)}.
\end{equation}
Pole pod krzywą ROC ($\mathrm{AUC}$) posiada jednoznaczną interpretację probabilistyczną: stanowi prawdopodobieństwo, że losowo wybrana obserwacja z klasy pozytywnej otrzyma od modelu wyższą ocenę punktową (score) niż losowo wybrana obserwacja z klasy negatywnej \cite{Fawcett2006,HanleyMcNeil1982}:
\begin{equation}
\mathrm{AUC} = P(\hat{y}(x^+) > \hat{y}(x^-)).
\end{equation}
Wartość $\mathrm{AUC}$ jest ściśle równoważna znormalizowanej statystyce Manna-Whitneya $U$ \cite{HanleyMcNeil1982}:
\begin{equation}
\mathrm{AUC} = \frac{U}{n^+ \cdot n^-} = \frac{R^+ - \frac{n^+(n^+ + 1)}{2}}{n^+ \cdot n^-},
\end{equation}
gdzie $R^+$ oznacza sumę rang przypisanych obserwacjom pozytywnym.

\subsection{Nieparametryczne testy istotności statystycznej}
W celu rzetelnego porównania metod czyszczenia danych zastosowano nieparametryczne testy rangowe:
\begin{enumerate}
    \item \textbf{Test rangowy znaków Wilcoxona} \cite{Wilcoxon1945}: nieparametryczny odpowiednik testu $t$-Studenta dla prób zależnych. Dla par wyników $(x_i, y_i)$ oblicza się różnice $d_i = x_i - y_i$, a następnie porządkuje ich wartości bezwzględne, przypisując im rangi $R_i$. Statystyka testowa wynosi:
    \begin{equation}
    W = \min(W^+, W^-), \quad W^+ = \sum_{d_i > 0} R_i, \quad W^- = \sum_{d_i < 0} R_i.
    \end{equation}
    \item \textbf{Test Friedmana} \cite{Friedman1937}: wielogrupowy test rangowy do weryfikacji hipotezy o braku różnic pomiędzy $k$ metodami ewaluowanymi na $N$ zbiorach prób. Statystyka testowa ma postać:
    \begin{equation}
    \chi_F^2 = \frac{12 N}{k(k + 1)} \left[ \sum_{j=1}^{k} \bar{R}_j^2 - \frac{k(k + 1)^2}{4} \right],
    \end{equation}
    gdzie $\bar{R}_j$ to średnia ranga $j$-tej metody.
\end{enumerate}
